\documentclass[10pt,a4paper]{amsart}
\usepackage[margin=19mm]{geometry}
\usepackage{amsmath,amssymb,mathtools}
\usepackage[scr=rsfs,cal=euler]{mathalfa}
\usepackage{xfrac}
\usepackage[unicode,hidelinks,bookmarks=false]{hyperref}
\newcommand{\ie}{i.e.\ }
\newcommand{\cf}{cf.\ }
\newcommand{\resp}{respectively\ }
\newcommand{\Subsection}{\subsection}
\providecommand{\nobreakdash}{\nobreak\hbox{-}}
\setlength{\emergencystretch}{2em}
\multlinegap0pt
\usepackage{amssymb}
\usepackage{algorithm}\usepackage{algpseudocode}
\usepackage{xcolor}
\usepackage{dsfont}
\usepackage{enumerate}
\usepackage{tikz}
\newtheorem{lemma}{Lemma}
\def\EE{\mathbb{E}}
\newcommand{\RR}{\mathbb{R}}
\newcommand{\Unif}{\mathrm{Uniform}}
\newcommand{\ZZ}{\mathbb{Z}}
\newcommand{\diag}{\mathop{\mathrm{diag}}}
\newcommand{\m}{\mathcal}
\title{High-Rate Quantized Matrix Multiplication II}
\author{Or Ordentlich, Yury Polyanskiy}
\date{Source: arXiv:2605.13768v2 (CC BY 4.0)}
\begin{document}
\maketitle

\noindent\emph{Source and licence.} High-Rate Quantized Matrix Multiplication II. Version arXiv:2605.13768v2 (CC BY 4.0), distributed by arXiv under the Creative Commons Attribution 4.0 International licence, http://creativecommons.org/licenses/by/4.0/, as recorded in the arXiv licence metadata for this exact version and shown on the abstract page https://arxiv.org/abs/2605.13768v2. Full licence notice: \texttt{licenses/arxiv-2605.13768-CC-BY-4.0.txt}.\par\smallskip

\noindent\textit{Editorial scope.} Counted unit: the article's lemma eq:DsicError (source0.tex lines 1130--1146). The article's complete original proof of it is delivered with it (source0.tex lines 1148--1159, one \texttt{proof} environment of 12 lines). No mathematics is written, repaired or re-derived: the statement and the proof are the article's own lines, in the article's own notation, and no retained line is substituted or re-flowed.  No further source-local context is needed: every cross-reference of every retained line resolves inside the two delivered spans.  Nothing was omitted from any retained span, so the package carries no omission record.  No retained line contains a citation, so the wrapper carries no bibliography and no bibliography entry.  No retained line contains a figure environment, an image-inclusion command or drawing code, so no image is delivered and the package declares no asset dependency.  Editorial formatting only, in addition: an AMS \texttt{amsart} preamble that restates the article's own declarations and macros used by the retained lines, namely the environments \texttt{lemma} on separate counters, together with the standard packages those lines need.  The retained environments print in this document as Lemma 1 in order of appearance, whereas the article prints the same items with the numbering of its own full document.  The complete native source of the article as distributed in the arXiv e-print for this exact version is bundled unchanged as \texttt{sources/200535/source0.tex}.
\begin{lemma}
Assume we apply the SIC Algorithm with $y\in\RR^n$, upper triangular $U\in\RR^{n\times n}$, and $\m{C}_i=\alpha_i\ZZ$ for all $i\in[n]$. Then
\begin{align}
 e_{\mathrm{SIC}}=y-U\cdot c_{\mathrm{SIC}} \in \m{P}_{\{U_{i,i},\alpha_i\}_{i=1}^n},\nonumber
\end{align}
where
\begin{align}
\m{P}_{\{U_{i,i},\alpha_i\}_{i=1}^n}=\prod_{i=1}^n \left[-\frac{|\alpha_i U_{i,i}|}{2},\frac{|\alpha_i U_{i,i}|}{2} \right). \nonumber
\end{align}
If in addition $e_{\mathrm{SIC}}\sim\Unif(\m{P}_{\{U_{i,i},\alpha_i\}_{i=1}^n})$ we have that
\begin{align}
D_{\mathrm{SIC}}&\triangleq\frac{1}{n}\EE\|Y-U\cdot c_{\mathrm{SIC}}\|^2\nonumber\\
&=\frac{1}{12}\cdot\frac{1}{n}\sum_{i=1}^n (\alpha_i U_{i,i})^2.  
\label{eq:DsicError}
\end{align}
\label{lem:FundCellSIC}
\end{lemma}

\begin{proof}
Let $c_{\mathrm{SIC}}(y,U)$ be the result of applying the SIC algorithm with inputs $y\in\RR^n$ and $U\in\RR^{n\times n}$. Denote $\m{A}=\diag(\alpha_1,\ldots,\alpha_n)\in\RR^{n\times n}$. The lemma follows from combining the two following observations:
\begin{enumerate}
\item $\m{P}_{\{U_{i,i},\alpha_i\}_{i=1}^n}=\left\{y\in\RR^n~:~c_{\mathrm{SIC}}(y,U)=0 \right\}$;
\item For any $z\in\ZZ^n$ it holds that $$c_{\mathrm{SIC}}(y+U \m{A} z,U)=\m{A}z+c_{\mathrm{SIC}}(y,U).$$
\end{enumerate}
Thus the SIC algorithm induces the partition of space to decision regions 
\begin{align}
\m{D}_{z}=U\m{A} z+\m{P}_{\{U_{i,i},\alpha_i\}_{i=1}^n},~\forall z\in\ZZ^n    
\end{align}
and $c_{\mathrm{SIC}}(y,U)=\m{A}z$ iff $y\in\m{D}_z$. Consequently, $e_{\mathrm{SIC}}=y-U\cdot c_{\mathrm{SIC}}(y,U)\in \m{P}_{\{U_{i,i},\alpha_i\}_{i=1}^n} $. The claim on $D_{\mathrm{SIC}}$ immediately follows from the product structure of $\m{P}_{\{U_{i,i},\alpha_i\}_{i=1}^n}$.
\end{proof}

\end{document}
